\uuid{AvSg}
\exo7id{3822}
\titre{exo7 3822}
\auteur{quercia}
\organisation{exo7}
\datecreate{2010-03-11}
\isIndication{false}
\isCorrection{true}
\chapitre{Espace euclidien, espace normé}
\sousChapitre{Problèmes matriciels}
\cpge{mpsi-mp2i.calcul-matriciel-systemes.e, pcsi.determinants.c, ptsi.determinants.a, mp-mpi.endomorphismes-euclidien.b, pc.endomorphismes-euclidiens.b, psi.endomorphismes-euclidien.b, pt.isometries-euclidien.b, bcpst2.produit-scalaire.c}
\module{Algèbre}
\niveau{L2}
\difficulte{}

\contenu{
\question{
Déterminer $\{M \in {\cal{M}}_n(\R) \text{ tq } M(^tMM)^2=I_n\}$.
}
\reponse{
$M = (^tMM)^{-2}$ est symétrique définie positive, donc diagonalisable
en base orthonormale. En examinant la forme diagonale on trouve $M=I$.
}
}
